\usepackage{amsmath}
\usepackage{xparse}
\usepackage{amsthm}
\usepackage{tikz}

% Seno
\def\sen{\mathop{\mbox{\normalfont sen}}\nolimits}

% Seno hiperbólico
\def\senh{\mathop{\mbox{\normalfont senh}}\nolimits}

% Arcoseno
\def\arcsen{\mathop{\mbox{\normalfont arcsen}}\nolimits}

% Gradiente
\def\grad{\mathop{\mbox{\normalfont grad}}\nolimits}

% Argumento principal
\def\Arg{\mathop{\mbox{\normalfont Arg}}\nolimits}

% Logaritmo complejo
\def\Log{\mathop{\mbox{\normalfont Log}}\nolimits}

% Rotacional
\def\rot{\mathop{\mbox{\normalfont rot}}\nolimits}

% Divergencia
\def\diver{\mathop{\mbox{\normalfont div}}\nolimits}

% Longitud
\def\longi{\mathop{\mbox{\normalfont long}}\nolimits}

% Índice de una curva
\def\ind{\mathop{\mbox{\normalfont Ind}}\nolimits}

% Rango de una matriz
\def\rang{\mathop{\mbox{\normalfont rang}}\nolimits}

% Norma
\newcommand{\norm}[1]{\lVert #1 \rVert}

% Derivada parcial
\newcommand{\derpar}[1]{\dfrac{\partial}{\partial #1}}

% Derivada parcial con función
\newcommand{\derparcial}[2]{\dfrac{\partial #1}{\partial #2}}

% Derivada parcial de orden n con funcion
\newcommand{\derparcialn}[3]{\dfrac{\partial^{#3} #1}{\partial #2^{#3}}}

% Integral doble en [a,b] x [c,d]
\newcommand{\intdob}[4]{\displaystyle \int_{#1}^{#2} \int_{#3}^{#4}}

% Integral triple en [a,b] x [c,d] x [e,f]
\newcommand{\inttri}[6]{\displaystyle \int_{#1}^{#2} \int_{#3}^{#4} \int_{#5}^{#6}}

% Números complejos
\def\C{\ensuremath{\mathbb{C}}}

% Números reales
\def\R{\ensuremath{\mathbb{R}}}

% Números racionales
\def\Q{\ensuremath{\mathbb{Q}}}

% Números enteros
\def\Z{\ensuremath{\mathbb{Z}}}

% Números naturales
\def\N{\ensuremath{\mathbb{N}}}

% Teorema
\newtheorem{teo}{Teorema}

% Corolario
\newtheorem{cor}{Corolario}

% Proposición
\newtheorem{prop}{Proposición}

% Lema
\newtheorem{lema}{Lema}

% Definición
\newtheorem{defi}{Definición}

% Observaciones
\newtheorem*{obs}{Observaciones}

\newtheorem*{ob}{Observación}

\newtheorem*{nota}{Nota}

\newtheorem*{problem}{Problema}

\newtheorem*{pregunta}{Pregunta}

\newtheorem{ejemplo}{Ejemplo}

\NewDocumentCommand{\overarrow}{O{=} O{\uparrow} m}{%
  \overset{\makebox[0pt]{\begin{tabular}{@{}c@{}}#3\\[0pt]\ensuremath{#2}\end{tabular}}}{#1}
}
\NewDocumentCommand{\underarrow}{O{=} O{\downarrow} m}{%
  \underset{\makebox[0pt]{\begin{tabular}{@{}c@{}}\ensuremath{#2}\\[0pt]#3\end{tabular}}}{\ensuremath{#1}}
}

\newcommand{\contradiction}{%
\begin{tikzpicture}[rotate=45,x=0.5ex,y=0.5ex]
\draw[color=red, line width=.1ex] (0,2) -- (3,2) (0,1) -- (3,1) (1,3) -- (1,0) (2,3) -- (2,0);
\end{tikzpicture}
}
